\documentclass{ip-lab}

\usepackage{url}
\usepackage{float}
\usepackage{graphicx}

\lablevel{n3}
\labnumber{Laboratorio 4}
\labtitle{Lectura de archivos y recorrido de diccionarios}

\makeheader

\begin{document}

\maketitle


\begin{sectionbox}{Preparación del ambiente de trabajo}
  Para este laboratorio se le proporcionará \textbf{los mismos insumos que para el proyecto del nivel 4}.
  
  Tendrá a su disposición un módulo de consola (\texttt{consola\_cupimundial.py}) que contiene un menú interactivo para probar sus funciones.

  Usted debe diligenciar el módulo de funciones llamado \texttt{cupimundial.py} en el que escribirá las funciones correspondientes a las actividades 1, 2 y 3 de este laboratorio. La consola importará automáticamente sus funciones desde este módulo.

  También se le proporcionará un archivo CSV llamado \texttt{cupimundial.csv} que contiene información sobre jugadores de las selecciones participantes en el torneo CupiMundial. El archivo CSV contiene datos organizados en columnas, donde cada fila representa un jugador diferente.
\end{sectionbox}

\pagebreak

\begin{sectionbox}{Estructura del archivo CSV}
  El archivo \texttt{cupimundial.csv} contiene información de jugadores organizada en las siguientes columnas:

  \begin{table}[H]
  \centering
  \caption{Descripción de las columnas del archivo: cupimundial.csv}
  \begin{tabular}{|c|l|p{4cm}|l|l|}
  \hline
  \textbf{Posición} & \textbf{Nombre de la columna} & \textbf{Descripción} & \textbf{Tipo} & \textbf{Ejemplo} \\
  \hline
  0 & \texttt{name} & Nombre del jugador & \texttt{str} & ``Lionel Messi'' \\
  \hline
  1 & \texttt{age} & Edad del jugador & \texttt{int} & 36 \\
  \hline
  2 & \texttt{team} & Nombre de la selección nacional a la que pertenece el jugador & \texttt{str} & ``Argentina'' \\
  \hline
  3 & \texttt{pace} & Velocidad del jugador (0 a 100) & \texttt{int} & 85 \\
  \hline
  4 & \texttt{shooting} & Habilidad de tiro del jugador (0 a 100) & \texttt{int} & 90 \\
  \hline
  5 & \texttt{passing} & Habilidad de pase del jugador (0 a 100) & \texttt{int} & 88 \\
  \hline
  6 & \texttt{dribbling} & Habilidad de regate del jugador (0 a 100) & \texttt{int} & 87 \\
  \hline
  7 & \texttt{position} & Posición del jugador. Valores posibles: ``Defender'', ``Forward'', ``Midfielder'', ``Reserve'', ``Striker'', ``Substitute'' o ``Winger'' & \texttt{str} & ``Winger'' \\
  \hline
  8 & \texttt{direction} & Dirección del jugador en el campo. Valores posibles: ``Left'', ``Center'', ``Right'' o ``Pending'' & \texttt{str} & ``Right'' \\
  \hline
  9 & \texttt{joined} & Fecha de ingreso del jugador al equipo en el formato ``YYYY-MM-DD'' & \texttt{str} & ``2023-01-01'' \\
  \hline
  10 & \texttt{contract} & Año de finalización del contrato en el formato ``YYYY'' & \texttt{str} & ``2026'' \\
  \hline
  11 & \texttt{preferred\_foot} & Pie preferido del jugador. Valores posibles: ``Right'' o ``Left'' & \texttt{str} & ``Right'' \\
  \hline
  12 & \texttt{international\_rep} & Reputación internacional del jugador (valor entero entre 1 y 5) & \texttt{int} & 5 \\
  \hline
  13 & \texttt{wage} & Salario mensual en USD del jugador & \texttt{int} & 320000 \\
  \hline
  \end{tabular}
  \end{table}
\end{sectionbox}

\begin{sectionbox}{Actividad 1: Cargar datos desde archivo CSV}
  Escriba una función \texttt{cargar\_cupimundial} que reciba como parámetro una cadena de caracteres con la ruta al archivo CSV y retorne un diccionario que organice la información de los jugadores según su selección nacional.

  Organice los datos en un diccionario, donde las llaves sean los nombres de las selecciones nacionales (strings no vacíos, sin espacios al inicio o al final). Los valores deben ser listas de diccionarios, donde cada elemento de la lista (cada diccionario) represente un jugador de esa selección, con sus atributos (a excepción del atributo \texttt{team}). Cada jugador se almacenará como un diccionario con todos sus campos, excepto el campo \texttt{team} (ya que la selección es la llave del diccionario principal). Los campos numéricos (\texttt{age}, \texttt{pace}, \texttt{shooting}, \texttt{passing}, \texttt{dribbling}, \texttt{international\_rep} y \texttt{wage}) deben convertirse a \texttt{int}.

  \begin{figure}[H]
    \centering
    \includegraphics[width=0.95\textwidth]{estructura.png}
    \caption{Ejemplo de la estructura deseada para el manejo de los datos.}
  \end{figure}

  \textbf{Estructura del retorno:}
  El diccionario retornado debe tener la siguiente estructura:
  \begin{verbatim}
  {
      "Argentina": [
          {
              "name": "Lionel Messi",
              "age": 36,
              "pace": 85,
              ...
          },
          ...
      ],
      "Poland": [
          {
              "name": "Robert Lewandowski",
              "age": 32,
              "pace": 78,
              ...
          },
          ...
      ],
      ...
  }
  \end{verbatim}
\end{sectionbox}

\begin{sectionbox}{Actividad 2: Buscar jugadores por posición y salario}
  Diligencia la función \texttt{buscar\_jugadores\_por\_posicion\_salario} que recibe como parámetros el diccionario de cupimundial (el mismo formato retornado por \texttt{cargar\_cupimundial}), una posición (string) y un salario máximo (entero), y retorna una lista con todos los jugadores que juegan en esa posición y cuyo salario sea menor o igual al salario máximo.
\end{sectionbox}

\begin{sectionbox}{Actividad 3: Recomendar jugador destacado}
  Diligencia la función \texttt{recomendar\_jugador\_destacado} que recibe como parámetros el diccionario de cupimundial, el nombre de una selección, una edad mínima, una edad máxima, una reputación internacional mínima y un pie preferido, y retorna el primer jugador que cumple simultáneamente con los siguientes criterios: juega en la posición más costosa (determinada por el salario total acumulado de quienes la ocupan), pertenece a la selección indicada, tiene una edad dentro del rango especificado, posee una reputación mayor o igual a la mínima y usa el pie preferido dado.

  \textbf{Pista:} La función \texttt{obtener\_posicion\_mas\_costosa} puede serle útil para implementar esta actividad.
\end{sectionbox}

\begin{sectionbox}{Entrega}
Entregue el archivo \texttt{cupimundial.py} a través de Bloque Neón en el laboratorio del Nivel 3 designado como ``N3-L4: Lectura de archivos y recorrido de diccionarios''.

Puede usar el módulo \texttt{consola\_cupimundial.py} proporcionado para probar sus funciones antes de la entrega. Su archivo debe funcionar correctamente desde la consola. Asegúrese de tener el archivo \texttt{cupimundial.csv} en el mismo directorio para probar su implementación.
\end{sectionbox}

\end{document}